\section{Finite observations and exact degree bookkeeping}
\label{sec:observations}

The probability measure on \(\{-1,1\}^N\) is uniform. For
\(S\subseteq[N]\), put \(\chi_S(x)=\prod_{i\in S}x_i\). Every real
function \(g\) has the expansion
\[
 g=\sum_{S\subseteq[N]}\widehat g(S)\chi_S,\qquad
 \widehat g(S)=\E[g\chi_S],
\]
and its Fourier degree is the largest \(|S|\) with
\(\widehat g(S)\ne0\). Nonzero constant functions have degree zero;
zero indicators can be discarded throughout.

\begin{definition}
A finite observation is a function \(F:\{-1,1\}^N\to\mathcal A\) with
finite range. Its cells are the nonempty sets \(C_a=\{F=a\}\). Define
\[
 D(F)=\degc(F)=\max_{\Pbb(F=a)>0}\deg\1_{C_a},\quad
 T_F=\E[H_N\mid F],\quad
 v(F)=\E T_F^2.
\]
The quantity \(N-D(F)\) is the \emph{Fourier codimension} of the
observation. It is a deficit of polynomial degree, not the geometric
codimension of its cells.
\end{definition}

Since \(\E H_N=0\), the score \(T_F\) is centered. Orthogonal projection,
or a direct cellwise expansion of the square, gives
\begin{equation}\label{eq:total-variance}
 N=v(F)+\E\Var(H_N\mid F).
\end{equation}
In particular \(0\le v(F)\le N\).

\begin{lemma}[Observation calculus]\label{lem:observation-calculus}
Every function of \(F\) has degree at most \(D(F)\). If \(F_i\) are
observations on disjoint blocks with degrees \(D_i\), then every
function of their tuple has degree at most \(\sum_iD_i\).
For the tuple itself, its exact cell degree is \(\sum_iD_i\), and its
retained variance is \(\sum_iv(F_i)\).
\end{lemma}

\begin{proof}
Expand a function of \(F\) as
\(\sum_a g(a)\1_{C_a}\). On disjoint blocks, a cell of the tuple is a
product of cell indicators, so its degree is at most the sum of the
degrees. Conversely, choose in each factor a cell of maximal degree
and a nonzero coefficient of that degree. Their product is a nonzero
coefficient on the disjoint union of the corresponding supports.
Hence the sum is attained. Independence gives
\[
 \E[H_{\mathrm{tot}}\mid(F_i)_i]=\sum_iT_{F_i}.
\]
The scores are independent and centered, proving variance additivity.
\end{proof}

If \(Q\) is any coarsening of the tuple, the tower property gives
\begin{equation}\label{eq:tower-scores}
 \E[H_{\mathrm{tot}}\mid Q]=
 \E\left[\sum_i T_{F_i}\,\middle|\,Q\right].
\end{equation}
Score collisions cause no degree difficulty: merging labels merely
adds their indicators.

\begin{lemma}[Boolean readout and the surviving linear bound]
\label{lem:readout}
Let \(v(F)>0\), put \(\kappa(F)=\E T_F^4/v(F)^2\), and define
\(f=\sgn(T_F)\), with \(\sgn(0)=1\). Then
\[
 \deg f\le D(F),\qquad
 L(f)=\E|T_F|\ge\sqrt{\frac{v(F)}{\kappa(F)}}.
\]
For every Boolean function \(g\),
\begin{equation}\label{eq:linear-bound}
 L_{\mathrm{abs}}(g)\le I(g)\le\deg g,
\end{equation}
where \(I(g)\) is total influence.
\end{lemma}

\begin{proof}
The degree assertion is Lemma~\ref{lem:observation-calculus}. Since
\(f\) is a function of \(F\),
\[
 L(f)=\E[fH_N]=\E[fT_F]=\E|T_F|.
\]
Hölder's inequality applied to
\(|T|^2=|T|^{2/3}|T|^{4/3}\) gives
\[
 \E T^2\le(\E|T|)^{2/3}(\E T^4)^{1/3},
\]
which rearranges to the desired bound.

For the last assertion define the discrete derivative
\[
 \partial_i g(x_{\ne i})
 =\frac{g(x_{\ne i},1)-g(x_{\ne i},-1)}2.
\]
It takes values in \(\{-1,0,1\}\), has mean
\(\widehat g(\{i\})\), and has second moment \(\Inf_i(g)\).
Thus \(|\widehat g(\{i\})|\le\Inf_i(g)\).
Orthogonality gives
\[
 I(g)=\sum_S|S|\widehat g(S)^2
 \le(\deg g)\sum_S\widehat g(S)^2=\deg g.
\]
This familiar inequality is also recorded in
\cite{NisanSzegedy1994,ODonnell2012}.
\end{proof}

A retained-variance excess \(v(F)>D(F)\) is therefore an intermediate
structural fact. It alone does not ensure that the sign readout
violates \(L(f)\le\sqrt{\deg f}\), because the first absolute moment
depends on the shape of the score distribution. Our twenty-bit example
makes this distinction concrete.
